\documentclass[10pt,a4paper]{amsart}
\usepackage{amsmath,amssymb,amsthm,graphicx}
% Retained preamble for the Stacks Project chapter morphisms.tex (Morphisms of Schemes).
% The chapter file pinned-original-source.tex carries only \input{preamble} on its first line, so no
% preamble text lives inside the chapter itself.  The Stacks Project uses one single shared preamble for
% every chapter; the snapshot stages it verbatim in sources/stacks-chow/preamble.tex.  The block below is
% copied mechanically, line for line, from that file's lines 14-18 (the Xy-pic setup) and lines 154-217
% (the native theorem declarations together with the mathematical macros).  Nothing was added, removed or
% redefined.  The wrapper class and the amsmath/amssymb/amsthm/graphicx base packages are supplied by the
% assembler, so the class/\externaldocument part of the shared preamble is not repeated here.
% For commutative diagrams we use Xy-pic
\usepackage[all]{xy}

% We use 2cell for 2-commutative diagrams.
\xyoption{2cell}
% Theorem environments.
%
\theoremstyle{plain}
\newtheorem{theorem}[subsection]{Theorem}
\newtheorem{proposition}[subsection]{Proposition}
\newtheorem{lemma}[subsection]{Lemma}

\theoremstyle{definition}
\newtheorem{definition}[subsection]{Definition}
\newtheorem{example}[subsection]{Example}
\newtheorem{exercise}[subsection]{Exercise}
\newtheorem{situation}[subsection]{Situation}

\theoremstyle{remark}
\newtheorem{remark}[subsection]{Remark}
\newtheorem{remarks}[subsection]{Remarks}

\numberwithin{equation}{subsection}

% Macros
%
\def\lim{\mathop{\mathrm{lim}}\nolimits}
\def\colim{\mathop{\mathrm{colim}}\nolimits}
\def\Spec{\mathop{\mathrm{Spec}}}
\def\Hom{\mathop{\mathrm{Hom}}\nolimits}
\def\Ext{\mathop{\mathrm{Ext}}\nolimits}
\def\SheafHom{\mathop{\mathcal{H}\!\mathit{om}}\nolimits}
\def\SheafExt{\mathop{\mathcal{E}\!\mathit{xt}}\nolimits}
\def\Sch{\mathit{Sch}}
\def\Mor{\mathop{\mathrm{Mor}}\nolimits}
\def\Ob{\mathop{\mathrm{Ob}}\nolimits}
\def\Sh{\mathop{\mathit{Sh}}\nolimits}
\def\NL{\mathop{N\!L}\nolimits}
\def\CH{\mathop{\mathrm{CH}}\nolimits}
\def\proetale{{pro\text{-}\acute{e}tale}}
\def\etale{{\acute{e}tale}}
\def\QCoh{\mathit{QCoh}}
\def\Ker{\mathop{\mathrm{Ker}}}
\def\Im{\mathop{\mathrm{Im}}}
\def\Coker{\mathop{\mathrm{Coker}}}
\def\Coim{\mathop{\mathrm{Coim}}}

% Boxtimes
%
\DeclareMathSymbol{\boxtimes}{\mathbin}{AMSa}{"02}

%
% Macros for moduli stacks/spaces
%
\def\QCohstack{\mathcal{QC}\!\mathit{oh}}
\def\Cohstack{\mathcal{C}\!\mathit{oh}}
\def\Spacesstack{\mathcal{S}\!\mathit{paces}}
\def\Quotfunctor{\mathrm{Quot}}
\def\Hilbfunctor{\mathrm{Hilb}}
\def\Curvesstack{\mathcal{C}\!\mathit{urves}}
\def\Polarizedstack{\mathcal{P}\!\mathit{olarized}}
\def\Complexesstack{\mathcal{C}\!\mathit{omplexes}}
% \Pic is the operator that assigns to X its picard group, usage \Pic(X)
% \Picardstack_{X/B} denotes the Picard stack of X over B
% \Picardfunctor_{X/B} denotes the Picard functor of X over B
\def\Pic{\mathop{\mathrm{Pic}}\nolimits}
\def\Picardstack{\mathcal{P}\!\mathit{ic}}
\def\Picardfunctor{\mathrm{Pic}}
\def\Deformationcategory{\mathcal{D}\!\mathit{ef}}

% --- appended 2026-10-12 (Chief): the project own definitions of the reference/slogan/history
% display environments, copied mechanically from the shared Stacks preamble lines 8-12
% (\usepackage{verbatim} plus the three \newenvironment lines, which the chapter preamble keeps
% verbatim from the project preamble). Nothing added, removed or redefined: these environments
% carry no mathematical content and the project itself typesets them as comments. Without these
% three lines a counted statement whose own text uses them fails with Environment slogan undefined.
\usepackage{verbatim}
\newenvironment{reference}{\comment}{\endcomment}
\newenvironment{slogan}{\comment}{\endcomment}
\newenvironment{history}{\comment}{\endcomment}
\usepackage{cleveref}
\crefname{originaldefinition}{original definition}{original definitions}
\Crefname{originaldefinition}{Original definition}{Original definitions}
\makeatletter
\newlabel{definition-property-local}{{01SR}{source}{}{sourcepointer-1}{}}
\newlabel{definition-property-local@cref}{{[originaldefinition][0][]01SR}{[1][0][]source}{}{}{}}
\makeatother
\crefname{originallemma}{original lemma}{original lemmas}
\Crefname{originallemma}{Original lemma}{Original lemmas}
\makeatletter
\newlabel{schemes-lemma-standard-open-two-affines}{{01IW}{source}{}{sourcepointer-2}{}}
\newlabel{schemes-lemma-standard-open-two-affines@cref}{{[originallemma][0][]01IW}{[1][0][]source}{}{}{}}
\makeatother
\title{A local property of ring maps holds on every pair of affine opens with the image of one inside the other}
\author{The Stacks Project Authors}
\date{stacks-project:morphisms.tex:local-snapshot-20261012; GNU FDL 1.2 or later}
\begin{document}
\maketitle
\noindent\textit{Editorial note.} All delivered mathematical text is the author's own, copied from the pinned Stacks Project chapter file pinned-original-source.tex; no statement or proof step was rewritten, reordered or removed.\par\medskip
\noindent\textit{Editorial note.} Following the multi-claim convention for this source, source\_url is the source-level URL https://stacks.math.columbia.edu/ and the source label lemma-locally-P is carried in claim\_key and here rather than in a per-claim tag URL.\par\medskip
\noindent\textit{Editorial note.} The preamble is the chapter's shared Stacks Project preamble, retained as sources/stacks-morphisms/retained-preamble.tex: a line-for-line excerpt of the shared Stacks preamble staged with this snapshot. No macro was invented or redefined, and this card adds no preamble text of its own.\par\medskip
\noindent\textit{Editorial note.} Statement span 2453-2462 and proof span 2464-2492 come from the halfA pool card operations/pool-cards/stacks-morphisms-halfA.json and were re-read in the pinned chapter file before this card was written.\par\medskip
\noindent\textit{Editorial note.} The proof refers to the three conditions of the cited definition by 'Definition (1)(a)', '(1)(b)' and '(1)(c)'; those references are the author's and are kept verbatim.\par\medskip
\noindent\textit{Editorial note.} 2 cross-references occur inside the counted spans; they are listed in external\_references as navigation pointers only.\par\medskip
\noindent\textit{Allowed standard prerequisites (editorial).}\begin{itemize}
\item The definition of a local property of ring maps with its conditions (a) on standard localizations of the target, (b) on standard localizations of the source and (c) on finite covers of the source. Cited repeatedly at proof lines 2687, 2696, 2699 and 2704. Every shrinking step in the proof is an application of one of these three conditions.
\item Two affine opens of a scheme with a common point have a common neighbourhood that is standard affine open in both, cited from the schemes chapter. Cited at proof lines 2684-2686 and 2691-2693. Produces the standard affine neighbourhoods used to invoke conditions (a) and (b) of locality.
\item An affine scheme is quasi-compact, so a cover of it by open pieces has a finite subcover. Used at proof lines 2700-2703. Turns the pointwise reductions into a finite cover to which condition (c) of locality can be applied.
\end{itemize}\medskip
\ifdefined\hypertarget\hypertarget{sourcepointer-1}{}\fi
\noindent\textit{Source reference.} Original definition 01SR: focused-seeds/AGP/sources/stacks-snapshot-20261011/tags.txt, source lines 2348--2348. Same-chapter reference used in the retained author text at source line(s) 2471, 2480, 2483, 2490, 2699, 3867, 3875, 4122, 5812, 7074, 7695, 8116. It is supplied only as a navigation pointer so that the retained text compiles; the counted item delivers the author's own statement and the author's complete proof, and no part of the referenced item is copied in or substituted for the counted argument.\par\medskip
\ifdefined\hypertarget\hypertarget{sourcepointer-2}{}\fi
\noindent\textit{Source reference.} Original lemma 01IW: focused-seeds/AGP/sources/stacks-snapshot-20261011/tags.txt, source lines 1993--1993. Same-chapter reference used in the retained author text at source line(s) 2470, 2476, 14942. It is supplied only as a navigation pointer so that the retained text compiles; the counted item delivers the author's own statement and the author's complete proof, and no part of the referenced item is copied in or substituted for the counted argument.\par\medskip
\begin{lemma}
\label{lemma-locally-P}
Let $f : X \to S$ be a morphism of schemes.
Let $P$ be a property of ring maps.
Let $U$ be an affine open of $X$,
and $V$ an affine open of $S$ such that
$f(U) \subset V$.
If $f$ is locally of type $P$ and $P$ is local,
then $P(\mathcal{O}_S(V) \to \mathcal{O}_X(U))$ holds.
\end{lemma}

\begin{proof}
As $f$ is locally of type $P$ for every $u \in U$ there exists an
affine open $U_u \subset X$ mapping into an affine open $V_u \subset S$
such that $P(\mathcal{O}_S(V_u) \to \mathcal{O}_X(U_u))$ holds.
Choose an open neighbourhood $U'_u \subset U \cap U_u$ of $u$
which is standard affine open in both $U$ and $U_u$, see
Schemes, Lemma \ref{schemes-lemma-standard-open-two-affines}.
By Definition \ref{definition-property-local} (1)(b)
we see that $P(\mathcal{O}_S(V_u) \to \mathcal{O}_X(U'_u))$ holds.
Hence we may assume that $U_u \subset U$ is a standard affine open.
Choose an open neighbourhood $V'_u \subset V \cap V_u$
of $f(u)$ which is standard affine open in both $V$ and $V_u$, see
Schemes, Lemma \ref{schemes-lemma-standard-open-two-affines}.
Then $U'_u = f^{-1}(V'_u) \cap U_u$ is a standard affine open
of $U_u$ (hence of $U$) and we have
$P(\mathcal{O}_S(V'_u) \to \mathcal{O}_X(U'_u))$ by
Definition \ref{definition-property-local} (1)(a).
Hence we may assume both $U_u \subset U$ and $V_u \subset V$
are standard affine open. Applying
Definition \ref{definition-property-local} (1)(b)
one more time we conclude that $P(\mathcal{O}_S(V) \to \mathcal{O}_X(U_u))$
holds. Because $U$ is quasi-compact we may choose a finite number
of points $u_1, \ldots, u_n \in U$ such that
$$
U = U_{u_1} \cup \ldots \cup U_{u_n}.
$$
By Definition \ref{definition-property-local} (1)(c)
we conclude that $P(\mathcal{O}_S(V) \to \mathcal{O}_X(U))$ holds.
\end{proof}

\end{document}
